\section{General-Case Compositions}\label{sec:general-case}

The synthetic fixtures are deliberate worst cases built to saturate the event cap.
To characterize behavior on realistic input, this section compiles four complete, idiomatic Motivo programs shipped with
Motivo Studio: \textit{Come As You Are}, a layered multi-track demo, a \textit{Pirates} theme, and \textit{F\"ur Elise}.
These exercise the language as a composer would, with patterns, parallel voices, multiple instrument tracks, chords, and
conditionals, rather than a single repeated statement.
Table~\ref{tab:general-case} reports their size and per-stage compile time.

\begin{table}[htbp]
  \centering
  \footnotesize
  \setlength{\tabcolsep}{4.5pt}
  \renewcommand{\arraystretch}{1.25}
  \begin{tabular}{lrrrrrrr}
    \toprule
    \textbf{Composition} & \textbf{Events} & \textbf{LOC} & \textbf{Parse} & \textbf{Semantic} & \textbf{Lowering} &
    \textbf{MIDI write} & \textbf{End-to-end} \\
    \midrule
    Come As You Are & 147 & 30  & 0.047 & 0.008 & 0.032 & 0.073 & $2.26 \pm 0.04$ \\
    Layered demo    & 188 & 101 & 0.120 & 0.021 & 0.061 & 0.082 & $2.43 \pm 0.09$ \\
    Pirates theme   & 191 & 62  & 0.124 & 0.023 & 0.084 & 0.078 & $2.42 \pm 0.06$ \\
    F\"ur Elise     & 294 & 58  & 0.090 & 0.019 & 0.145 & 0.088 & $2.47 \pm 0.06$ \\
    \bottomrule
  \end{tabular}
  \caption{Per-stage mean compile time (ms) for four real compositions on the current \texttt{motivoc}; the end-to-end
    column reports the mean $\pm$ standard deviation over 50 runs. The four stage timers sum to $0.15$--$0.35$~ms; the
  gap to the end-to-end column is process startup and file I/O.}\label{tab:general-case}
\end{table}

Two points stand out.
First, the realistic expansion ratio is modest: these pieces emit between $1.9$ and $5.1$ events per source line, two to
three orders of magnitude below the synthetic fixtures, which reach roughly $2{,}200$ events per line at the cap.
Real music interleaves largely unique material, so the dramatic compression of Section~\ref{sec:compression} is a
property of repetitive structure, not of typical compositions; the honest expectation for hand-written pieces is a small
constant factor.
This tempers the synthetic compression figures rather than contradicting them: the language still removes manual event
placement, but the leverage depends on how much genuine repetition a piece contains.
The relationship between length and stages is visible in the table: the most verbose sources (the 101-line demo and the
Pirates theme) have the highest parse times, while the piece with the most events (\textit{F\"ur Elise}) has the highest
lowering and MIDI time.

Second, every stage is fast.
The four pipeline stages together account for only $0.15$--$0.35$~ms; the remaining $2$~ms or so of each compile is the
fixed process-startup floor identified in Section~\ref{detail-subsec:scaling-sweep}.
For compositions of this scale, compile time is effectively constant and well below perceptible latency, which is what
makes the studio's compile-to-preview loop feel immediate
(Section~\ref{detail-subsec:execution-speed-and-real-time-viability}).
The synthetic fixtures show where cost eventually appears; the general case shows that realistic programs never approach
it.
